Die Brillengläser haben draussen die tiefe Temperatur der Aussenluft angenommen. Die warme Raumluft enthält Wasserdampf. Die Luftschicht unmittelbar an den kalten Gläsern kühlt sich stark ab, die absolute Luftfeuchtigkeit bleibt dabei gleich. Weil die Sättigungsdichte des Wasserdampfs mit sinkender Temperatur stark abnimmt, würde die relative Luftfeuchtigkeit in dieser Schicht auf über hundert Prozent steigen: \result{Die Luft wird unter ihren Taupunkt abgekühlt, und der überschüssige Wasserdampf kondensiert als feine Tröpfchen auf dem Glas.}

Sobald sich die Gläser über den Taupunkt der Raumluft erwärmt haben, verdunsten die Tröpfchen wieder, und die Brille wird klar.
